% https://www.overleaf.com/latex/examples/latex-portrait-poster-template/gybjbztdkvyg
%\title{LaTeX Portrait Poster Template}
%%%%%%%%%%%%%%%%%%%%%%%%%%%%%%%%%%%%%%%%%
% a0poster Portrait Poster
% LaTeX Template
% Version 1.0 (22/06/13)
%
% The a0poster class was created by:
% Gerlinde Kettl and Matthias Weiser (tex@kettl.de)
% 
% This template has been downloaded from:
% http://www.LaTeXTemplates.com
%
% License:
% CC BY-NC-SA 3.0 (http://creativecommons.org/licenses/by-nc-sa/3.0/)
%
%%%%%%%%%%%%%%%%%%%%%%%%%%%%%%%%%%%%%%%%%

%----------------------------------------------------------------------------------------
%	PACKAGES AND OTHER DOCUMENT CONFIGURATIONS
%----------------------------------------------------------------------------------------

\documentclass[a0,portrait]{a0poster}
\usepackage{fontspec}

\usepackage{multicol} % This is so we can have multiple columns of text side-by-side
\columnsep=100pt % This is the amount of white space between the columns in the poster
\columnseprule=3pt % This is the thickness of the black line between the columns in the poster

\usepackage[svgnames]{xcolor} % Specify colors by their 'svgnames', for a full list of all colors available see here: http://www.latextemplates.com/svgnames-colors

% \usepackage{times} % Use the times font
%\usepackage{palatino} % Uncomment to use the Palatino font

\usepackage{graphicx} % Required for including images
\graphicspath{{img/}} % Location of the graphics files
\usepackage{booktabs} % Top and bottom rules for table
\usepackage[font=small,labelfont=bf]{caption} % Required for specifying captions to tables and figures
% \usepackage{amsfonts, amsmath, amsthm, amssymb} % For math fonts, symbols and environments
\usepackage{wrapfig} % Allows wrapping text around tables and figures

\usepackage{float}

%\newfontfamily{\defaultfont}{Linux Libertine O}
\newfontfamily{\defaultfont}{TeX Gyre Schola}

\newcommand{\uname}[1]{{\index{#1}'\defaultfont\textsc{\MakeLowercase{#1}}'}}
\newcommand{\ucode}[1]{{\defaultfont\texttt{U+#1}}}
\newcommand{\mname}[1]{\index{#1}'MUFI \textsc{\MakeLowercase{#1}}'}
\newcommand{\mcode}[1]{\texttt{M+#1}}
\newcommand{\aname}[1]{\textsc{#1}}
\newcommand{\acode}[1]{\textsc{#1}}
\newcommand{\prname}[1]{\textsf{#1}}

\newfontfamily{\mufi}{Junicode}
\newcommand{\mufitags}[1]{{\mufi\addfontfeature{StylisticSet=10}#1}}

\usepackage{hyperref}

\begin{document}
\defaultfont

%----------------------------------------------------------------------------------------
%	POSTER HEADER 
%----------------------------------------------------------------------------------------

% The header is divided into two boxes:
% The first is 75% wide and houses the title, subtitle, names, university/organization and contact information
% The second is 25% wide and houses a logo for your university/organization or a photo of you
% The widths of these boxes can be easily edited to accommodate your content as you see fit

\begin{minipage}[b]{0.75\linewidth}
\veryHuge \color{NavyBlue} \textbf{Developing a tool for glyph spotting\\ in scanned printed texts} \color{Black}\\[0.25cm] % Title
% \Huge\textit{a
%   computer resource}\\[2cm] % Subtitle
\huge \textbf{Janusz S. Bień} (ORCID 0000-0001-5006-8183) \\[0.5cm] % Author(s)
\huge University of Warsaw (retired)\\[0.4cm] % University/organization
\Large \texttt{jsbien@uw.edu.pl}, \url{https://jsbien.github.io/}\\
% telefon --- 1 (000) 111 1111\\
\end{minipage}
%
\begin{minipage}[b]{0.25\linewidth}
\includegraphics[width=10cm]{img/jsbien.png}\\
\end{minipage}

\vspace{1cm} % A bit of extra whitespace between the header and poster content

%----------------------------------------------------------------------------------------

\begin{multicols}{2} % This is how many columns your poster will be broken into, a portrait poster is generally split into 2 columns

%----------------------------------------------------------------------------------------
%	ABSTRACT
%----------------------------------------------------------------------------------------

\color{Navy} % Navy color for the abstract

\begin{abstract}
  In document-image retrieval, word spotting (keyword spotting, KWS)
  finds occurrences of a query image in page images without full OCR,
  typically by retrieving similar word images using segmentation-based
  candidates or segmentation-free sliding-window search. Many modern
  KWS systems map each word image to a learned feature vector so that
  instances of the same word lie close together, enabling efficient
  nearest-neighbor retrieval. Our tool is related in spirit but
  targets glyph-level retrieval in historical printed material by
  exploiting a DjVu-specific property: shape dictionaries with
  explicit occurrence lists produced by the DjVu encoding
  pipeline for compression purposes.
  % This shifts emphasis from candidate generation to
  % interactive inspection, filtering, and curation of automatically
  % clustered shape classes, supporting a curated glyph image collection
  % (“glyph gallery”) and typographic inventories when
  % segmentation—often the bottleneck in query-by-example KWS—is
  % unreliable.
\end{abstract}

%----------------------------------------------------------------------------------------
%	INTRODUCTION
%----------------------------------------------------------------------------------------

\color{SaddleBrown} % SaddleBrown color for the introduction

\section*{Introduction}

The DjVu format was once popular in  digital libraries, including
the Internet Archive.\footnote{\url{https://en.wikipedia.org/wiki/DjVu}}
However, not every DjVu document contains the shape dictionaries we are
interested in. If not, it has to be regenerated, e.g. using Jakub Wilk's
\textsf{didjvu} tool, currently maintained by Friedrich
Froebel\footnote{\url{https://github.com/FriedrichFroebel/didjvu}}.
% Fröbel

A browser for the shape dictionary was developed by Michał Rudolf
(later, a major contribution was made by Alexander
Trufanov).\footnote{\url{https://github.com/jsbien/djview4shapes}}
Unfortunately the tool, see Fig.~\ref{fig:K-C-e}, in practice appeared
not sufficiently convenient.

%\begin{figure}[ht]
\begin{figure}[H]
  \centering
  \includegraphics[width=\hsize]{lB42}
  \caption{\prname{djview4shapes} and Gutenberg's bible}
  \label{fig:K-C-e}
\end{figure}

I also designed a more sophisticated tool and an attempt to implement
it was made by Piotr Sikora. The idea was to have a database of shapes
from many documents that could be updated and queried remotely, and a
client used to browse and label the
shapes.\footnote{\url{https://github.com/jsbien/ndt/wiki/z_shapes}}
Unfortunately the tool was prohibitively slow and there was neither
the opportunity nor motivation to try to fix it.

%\begin{figure}[ht]
\begin{figure}[H]
  \centering
  \includegraphics[width=\hsize]{e1_01.png}
  \caption{Piotr Sikora's shape labelling client}
  \label{fig:Sikora}
\end{figure}

Recently I had a fresh look at Sikora's \textit{exportshapes} utility
used for exporting shapes to a SQL compatible database and found that
it is still usable and serves well its purpose, the problem is only
with the client. However in the meantime my needs changed and I
decided to design and develop a completely new client. I am not
a programmer but the availability of AI chat tools made this task feasible.
I present here a progress report.
  % ----------------------------------------------------------------------------------------

\color{DarkSlateGray} % DarkSlateGray color for the rest of the content



%\small

\section{The repository}

The tool is developed in a standard public \textsf{Git} repository on
GitHub, the software is available under an open-source
license.\footnote{\url{https://github.com/jsbien/shape-browser4djvu}}. The
repository site has standard tabs for reporting issues and for
discussions. Wiki is enabled but not used at present.

The repository has several so-called branches. Besides those used for
the development (the good practice is to test first the new features
in a separate branch) an important one is the Sikora branch with the
source of the \textsf{exportshapes} tool.

The tool is developed in Python on a Debian GNU/Linux
platform. Porting to Windows should be possible, but I'm not
interested in it.

At this stage, GUI aesthetics are a low priority.

For testing purposes I use \textit{[Exercitium veteris
  artis]. Argumentum in librum Porphyrii} published in 1504 by Jan
Haller, the digitization is available in Wielkopolska Digital
Library.\footnote{\url{https://www.wbc.poznan.pl/dlibra/publication/559203/edition/479475}}
Its DjVu version has 586\,035 shapes.
\section{The shape browser}

Figure~\ref{fig:djview} illustrates some of the features of the tool.

It should be noted that the shapes in the shape dictionaries are
organized as trees. It should be also noted that the shapes are just
connected components (i.e. groups of neighboring black pixels in the
binary mask), so e.g.  diacritics are separated from the base letter,
on the other hand characters merged by ink bleed are considered to be
a single shape. Despite these drawbacks I consider the shape
dictionary a good first approximation of the character inventory for a
given text.

In the figure, one of the root shapes is selected (in the red frame) and
in the  bottom-right corner we see a panel with its properties. From
a practical standpoint, the most important information is about the
occurrences of the shape in the text: there is one occurrence on page 111
at the coordinates x=780 and y=645. Clicking on this information
opens the \textsf{djview} viewer with the occurrence
centered on the screen. If we want to see a larger context, we can
simply use the \textsf{djview} zoom commands.

\begin{figure}[H]
  \centering
  \includegraphics[width=\hsize]{djview}
  \caption{A sample shape}
  \label{fig:djview}
\end{figure}

\begin{figure}[H]
  \centering
  \includegraphics[width=\hsize]{djview_zoom}
  \caption{A shape context: \textit{Arguitur quarto:
      \textbf{substantia} non est genus
      \ldots}\protect\footnotemark}
  \label{fig:djview_zoom}
\end{figure}
\footnotetext{\url{https://latin.stackexchange.com/questions/27099/a-brevigraph-meaning}}
% https://latin.stackexchange.com/questions/27099/a-brevigraph-meaning
% It's not my area but at a guess I would say substantia. – Cairnarvon
% Cairnarvon may be right, but I don't recognise this as a common
% abbreviation, except that b followed by a curl like that would
% normally be something like ber. Before that it is indeed Arguituir
% quarto, and after non est genus. Can you post more context, a much
% larger part of the text? The link doesn't work for me. – Cerberus


In the top-left corner we have several bars: menu bar, banner bar with
some information about the current display, and the most important
filter bar. As the number of shapes can be very large, we can limit
the number of displayed shapes in several ways. In particular it is
practical to browse first just the roots of the shape trees.



\section{Final remark}
\label{sec:final-remarks}
I'm developing the tool for my personal use without any funding or
institutional support.  I hope to use it to supplement some items in
the so-called Typoglyph Gallery with the Unicode encoding of the
glyphs which sometimes are far from
obvious\footnote{\url{https://jsbien.github.io/typoglyphs/}}. I would
be glad if the tool also proved useful for others.

The text was proofread by ChatGPT, which also drafted the abstract.
\end{multicols}
\end{document}

%%% Local Variables:
%%% mode: latex
%%% ispell-local-dictionary: "english"
%%% TeX-master: t
%%% coding: utf-8-unix
%%% TeX-PDF-mode: t
%%% TeX-engine: xetex
%%% End:

